% Options for packages loaded elsewhere
% Options for packages loaded elsewhere
\PassOptionsToPackage{unicode}{hyperref}
\PassOptionsToPackage{hyphens}{url}
\PassOptionsToPackage{dvipsnames,svgnames,x11names}{xcolor}
%
\documentclass[
  11pt,
]{article}
\usepackage{xcolor}
\usepackage[margin=1in]{geometry}
\usepackage{amsmath,amssymb}
\setcounter{secnumdepth}{-\maxdimen} % remove section numbering
\usepackage{iftex}
\ifPDFTeX
  \usepackage[T1]{fontenc}
  \usepackage[utf8]{inputenc}
  \usepackage{textcomp} % provide euro and other symbols
\else % if luatex or xetex
  \usepackage{unicode-math} % this also loads fontspec
  \defaultfontfeatures{Scale=MatchLowercase}
  \defaultfontfeatures[\rmfamily]{Ligatures=TeX,Scale=1}
\fi
\usepackage{lmodern}
\ifPDFTeX\else
  % xetex/luatex font selection
\fi
% Use upquote if available, for straight quotes in verbatim environments
\IfFileExists{upquote.sty}{\usepackage{upquote}}{}
\IfFileExists{microtype.sty}{% use microtype if available
  \usepackage[]{microtype}
  \UseMicrotypeSet[protrusion]{basicmath} % disable protrusion for tt fonts
}{}
\makeatletter
\@ifundefined{KOMAClassName}{% if non-KOMA class
  \IfFileExists{parskip.sty}{%
    \usepackage{parskip}
  }{% else
    \setlength{\parindent}{0pt}
    \setlength{\parskip}{6pt plus 2pt minus 1pt}}
}{% if KOMA class
  \KOMAoptions{parskip=half}}
\makeatother
% Make \paragraph and \subparagraph free-standing
\makeatletter
\ifx\paragraph\undefined\else
  \let\oldparagraph\paragraph
  \renewcommand{\paragraph}{
    \@ifstar
      \xxxParagraphStar
      \xxxParagraphNoStar
  }
  \newcommand{\xxxParagraphStar}[1]{\oldparagraph*{#1}\mbox{}}
  \newcommand{\xxxParagraphNoStar}[1]{\oldparagraph{#1}\mbox{}}
\fi
\ifx\subparagraph\undefined\else
  \let\oldsubparagraph\subparagraph
  \renewcommand{\subparagraph}{
    \@ifstar
      \xxxSubParagraphStar
      \xxxSubParagraphNoStar
  }
  \newcommand{\xxxSubParagraphStar}[1]{\oldsubparagraph*{#1}\mbox{}}
  \newcommand{\xxxSubParagraphNoStar}[1]{\oldsubparagraph{#1}\mbox{}}
\fi
\makeatother


\usepackage{longtable,booktabs,array}
\newcounter{none} % for unnumbered tables
\usepackage{calc} % for calculating minipage widths
% Correct order of tables after \paragraph or \subparagraph
\usepackage{etoolbox}
\makeatletter
\patchcmd\longtable{\par}{\if@noskipsec\mbox{}\fi\par}{}{}
\makeatother
% Allow footnotes in longtable head/foot
\IfFileExists{footnotehyper.sty}{\usepackage{footnotehyper}}{\usepackage{footnote}}
\makesavenoteenv{longtable}
\usepackage{graphicx}
\makeatletter
\newsavebox\pandoc@box
\newcommand*\pandocbounded[1]{% scales image to fit in text height/width
  \sbox\pandoc@box{#1}%
  \Gscale@div\@tempa{\textheight}{\dimexpr\ht\pandoc@box+\dp\pandoc@box\relax}%
  \Gscale@div\@tempb{\linewidth}{\wd\pandoc@box}%
  \ifdim\@tempb\p@<\@tempa\p@\let\@tempa\@tempb\fi% select the smaller of both
  \ifdim\@tempa\p@<\p@\scalebox{\@tempa}{\usebox\pandoc@box}%
  \else\usebox{\pandoc@box}%
  \fi%
}
% Set default figure placement to htbp
\def\fps@figure{htbp}
\makeatother





\setlength{\emergencystretch}{3em} % prevent overfull lines

\providecommand{\tightlist}{%
  \setlength{\itemsep}{0pt}\setlength{\parskip}{0pt}}



 


\usepackage{amsmath,amssymb,enumitem,booktabs}
\setlist{itemsep=2pt,topsep=3pt}
\usepackage{fancyhdr}\pagestyle{fancy}\fancyhf{}
\fancyhead[L]{\small DS701 Fall 2026}\fancyhead[R]{\small Homework 5}
\fancyfoot[C]{\small\thepage}
\renewcommand{\headrulewidth}{0.4pt}
\usepackage[breakable]{tcolorbox}
% --- answer-space macros (problems PDF only; no-ops in the solutions PDF) ---
% \answerspace{2in}  vertical room for handwritten work, sized from the solution
% \blank             a short fill-in rule; \blank[1.5in] for a wider one
% \answerpage        start the next problem on a fresh page (Gradescope page matching)
\newcommand{\answerspace}[1]{\par\vspace*{#1}}  % \vspace* so room is not dropped at a page break
\newcommand{\blank}[1][1.2in]{\rule{#1}{0.4pt}}
\newcommand{\answerpage}{\clearpage}
\newtcolorbox{solutionbox}{breakable,colback=blue!4,colframe=blue!45!black,title=Solution,fonttitle=\bfseries,left=4pt,right=4pt,top=3pt,bottom=3pt}
\makeatletter
\@ifpackageloaded{caption}{}{\usepackage{caption}}
\AtBeginDocument{%
\ifdefined\contentsname
  \renewcommand*\contentsname{Table of contents}
\else
  \newcommand\contentsname{Table of contents}
\fi
\ifdefined\listfigurename
  \renewcommand*\listfigurename{List of Figures}
\else
  \newcommand\listfigurename{List of Figures}
\fi
\ifdefined\listtablename
  \renewcommand*\listtablename{List of Tables}
\else
  \newcommand\listtablename{List of Tables}
\fi
\ifdefined\figurename
  \renewcommand*\figurename{Figure}
\else
  \newcommand\figurename{Figure}
\fi
\ifdefined\tablename
  \renewcommand*\tablename{Table}
\else
  \newcommand\tablename{Table}
\fi
}
\@ifpackageloaded{float}{}{\usepackage{float}}
\floatstyle{ruled}
\@ifundefined{c@chapter}{\newfloat{codelisting}{h}{lop}}{\newfloat{codelisting}{h}{lop}[chapter]}
\floatname{codelisting}{Listing}
\newcommand*\listoflistings{\listof{codelisting}{List of Listings}}
\makeatother
\makeatletter
\makeatother
\makeatletter
\@ifpackageloaded{caption}{}{\usepackage{caption}}
\@ifpackageloaded{subcaption}{}{\usepackage{subcaption}}
\makeatother
\usepackage{bookmark}
\IfFileExists{xurl.sty}{\usepackage{xurl}}{} % add URL line breaks if available
\urlstyle{same}
% fallback for those not using the hyperref driver hyperxmp:
\makeatletter
\@ifundefined{xmpquote}{\newcommand{\xmpquote}[1]{#1}}{}
\makeatother
\hypersetup{
  pdftitle={DS701 Homework 5: Bias--Variance{,} Trees{,} and Ensembles},
  colorlinks=true,
  linkcolor={blue},
  filecolor={Maroon},
  citecolor={Blue},
  urlcolor={Blue},
  pdfcreator={LaTeX via pandoc}}


\title{DS701 Homework 5: Bias--Variance, Trees, and Ensembles}
\usepackage{etoolbox}
\makeatletter
\providecommand{\subtitle}[1]{% add subtitle to \maketitle
  \apptocmd{\@title}{\par {\large #1 \par}}{}{}
}
\makeatother
\subtitle{Released Wed Oct 7 · due Tue Oct 13, 11:59 pm on Gradescope}
\author{}
\date{}
\begin{document}
\maketitle


\noindent\textbf{Name:}\ \rule{2.0in}{0.4pt}\qquad\textbf{BU ID:}\ \rule{1.2in}{0.4pt}
\vspace{6pt}

\textbf{How this works.} Answer \textbf{by hand on paper} (or a tablet
with a stylus), show your work, then scan and upload to Gradescope,
matching each problem to its pages. This is deliberate practice for the
closed-notes quizzes, which are in the same style --- so do it as much
as possible without a computer or an AI assistant. If you are stuck, you
can try AI in tutorial style or reach out to course staff. It is
\textbf{participation-graded}: a serious attempt at every problem earns
full credit, and worked solutions are released the day after the
deadline. Budget 1--2 hours.

\textbf{Covers:} Lecture 07 (Generalization) and Lecture 08 (Decision
Trees and Random Forests).

\subsection{Problem 1: Bias and variance, made
concrete}\label{problem-1-bias-and-variance-made-concrete}

\emph{Practices: turning ``high bias / high variance'' from words into
numbers, and reading a learning curve.}

We want to predict the value of some function at a single test input
\(x_0\), whose true value is \(f(x_0) = 10\). Two models were each
trained four times, on four independent training sets drawn from the
same distribution, and asked to predict at \(x_0\):

{\def\LTcaptype{none} % do not increment counter
\begin{longtable}[]{@{}lllll@{}}
\toprule\noalign{}
training set & 1 & 2 & 3 & 4 \\
\midrule\noalign{}
\endhead
\bottomrule\noalign{}
\endlastfoot
Model A (very simple) & 7 & 8 & 8 & 9 \\
Model B (very flexible) & 6 & 14 & 8 & 12 \\
\end{longtable}
}

For a model whose predictions at \(x_0\) across training sets are
\(\hat y_1,\dots,\hat y_4\) with mean \(\bar y\), use these definitions:
\[\text{bias} = \bar y - f(x_0),\qquad
\text{variance} = \tfrac14\textstyle\sum_{i=1}^{4}(\hat y_i-\bar y)^2,\qquad
\text{expected squared error} \approx \text{bias}^2 + \text{variance} + \sigma^2,\]
where \(\sigma^2\) is the variance of the label noise (irreducible; the
same for every model).

\newpage

\begin{enumerate}
\def\labelenumi{(\alph{enumi})}
\tightlist
\item
  For each model compute the mean prediction, bias\(^2\), and variance.
  Which model is ``high bias'' and which is ``high variance''? Note that
  each is \emph{low} in the other quantity.
\end{enumerate}

{\def\LTcaptype{none} % do not increment counter
\begin{longtable}[]{@{}
  >{\raggedright\arraybackslash}p{(\linewidth - 8\tabcolsep) * \real{0.2000}}
  >{\raggedright\arraybackslash}p{(\linewidth - 8\tabcolsep) * \real{0.2000}}
  >{\raggedright\arraybackslash}p{(\linewidth - 8\tabcolsep) * \real{0.2000}}
  >{\raggedright\arraybackslash}p{(\linewidth - 8\tabcolsep) * \real{0.2000}}
  >{\raggedright\arraybackslash}p{(\linewidth - 8\tabcolsep) * \real{0.2000}}@{}}
\toprule\noalign{}
\begin{minipage}[b]{\linewidth}\raggedright
\end{minipage} & \begin{minipage}[b]{\linewidth}\raggedright
mean \(\bar y\)
\end{minipage} & \begin{minipage}[b]{\linewidth}\raggedright
bias\(^2\)
\end{minipage} & \begin{minipage}[b]{\linewidth}\raggedright
variance
\end{minipage} & \begin{minipage}[b]{\linewidth}\raggedright
high bias or high variance?
\end{minipage} \\
\midrule\noalign{}
\endhead
\bottomrule\noalign{}
\endlastfoot
Model A & \blank[0.7in] & \blank[0.7in] & \blank[0.7in] &
\blank[1.4in] \\
Model B & \blank[0.7in] & \blank[0.7in] & \blank[0.7in] &
\blank[1.4in] \\
\end{longtable}
}

\answerspace{2.5in}

\begin{enumerate}
\def\labelenumi{(\alph{enumi})}
\setcounter{enumi}{1}
\tightlist
\item
  With \(\sigma^2 = 3\), compute the expected squared error of each
  model at \(x_0\) and state which model you would deploy. Would your
  choice change for any value of \(\sigma^2\)? One sentence.
\end{enumerate}

Expected squared error --- A: \blank[0.8in]\quad B:
\blank[0.8in]\quad Deploy: \blank[0.8in]

\answerspace{1.5in}
\answerpage

\begin{enumerate}
\def\labelenumi{(\alph{enumi})}
\setcounter{enumi}{2}
\tightlist
\item
  The table below is a \emph{learning curve} --- mean squared error on
  the training set and on a held-out validation set as the training-set
  size \(n\) grows --- for two other models, X and Y (all four columns
  come from the same data-generating process).
\end{enumerate}

{\def\LTcaptype{none} % do not increment counter
\begin{longtable}[]{@{}lllll@{}}
\toprule\noalign{}
\(n\) & X: train & X: validation & Y: train & Y: validation \\
\midrule\noalign{}
\endhead
\bottomrule\noalign{}
\endlastfoot
50 & 3.9 & 8.0 & 0.1 & 9.0 \\
200 & 4.0 & 5.5 & 0.2 & 6.5 \\
1000 & 4.1 & 4.5 & 0.3 & 4.5 \\
5000 & 4.1 & 4.2 & 0.4 & 3.0 \\
\end{longtable}
}

Diagnose each of X and Y as mainly high-bias or mainly high-variance,
giving the feature of the curve that tells you. For each, say whether
collecting five times more data is worth the money, and name one other
remedy from Lecture 07.

X is mainly high \blank[1in]\quad Y is mainly high \blank[1in]

\answerspace{2.5in}

\answerpage

\subsection{Problem 2: Choosing a split, then knowing when to
stop}\label{problem-2-choosing-a-split-then-knowing-when-to-stop}

\emph{Practices: entropy and Gini by hand, information gain, and why the
greedy tree keeps growing.}

A lender records, for ten past borrowers, whether they own their home,
whether they are married, and whether they later \textbf{defaulted} (the
label):

{\def\LTcaptype{none} % do not increment counter
\begin{longtable}[]{@{}llll@{}}
\toprule\noalign{}
\# & Home owner & Married & Default \\
\midrule\noalign{}
\endhead
\bottomrule\noalign{}
\endlastfoot
1 & Y & Y & N \\
2 & Y & Y & N \\
3 & Y & N & N \\
4 & Y & N & N \\
5 & N & Y & Y \\
6 & N & Y & N \\
7 & N & Y & N \\
8 & N & N & Y \\
9 & N & N & Y \\
10 & N & N & Y \\
\end{longtable}
}

With \(p_i\) the fraction of class \(i\) at a node (and
\(0\log_2 0 = 0\)), use
\[\text{Entropy}=-\sum_i p_i\log_2 p_i,\qquad \text{Gini}=1-\sum_i p_i^2 .\]
The collective impurity of \(k\) children is their size-weighted average
\(I(\text{children})=\sum_{j=1}^{k}\frac{N(v_j)}{N}\,I(v_j)\), and the
information gain of a split is \(I(\text{parent})-I(\text{children})\).
You may use \(\log_2 3 \approx 1.585\) and \(\log_2 5 \approx 2.322\);
every other logarithm you need follows from these
(e.g.~\(\log_2\tfrac{2}{5} = 1 - \log_2 5\)). Round to 3 decimals.

\begin{enumerate}
\def\labelenumi{(\alph{enumi})}
\tightlist
\item
  Compute the entropy and the Gini index of the root node (all ten
  rows).
\end{enumerate}

Entropy: \blank[0.8in]\quad Gini: \blank[0.8in]

\answerspace{1.8in}
\answerpage

\begin{enumerate}
\def\labelenumi{(\alph{enumi})}
\setcounter{enumi}{1}
\tightlist
\item
  Two candidate splits are available at the root: on \emph{Home owner}
  and on \emph{Married}. For each, compute the collective (weighted)
  entropy of the two children and the information gain. Which split does
  the greedy algorithm choose? Check that the Gini gains agree on the
  ordering (compute both weighted Gini values).
\end{enumerate}

{\def\LTcaptype{none} % do not increment counter
\begin{longtable}[]{@{}
  >{\raggedright\arraybackslash}p{(\linewidth - 8\tabcolsep) * \real{0.2000}}
  >{\raggedright\arraybackslash}p{(\linewidth - 8\tabcolsep) * \real{0.2000}}
  >{\raggedright\arraybackslash}p{(\linewidth - 8\tabcolsep) * \real{0.2000}}
  >{\raggedright\arraybackslash}p{(\linewidth - 8\tabcolsep) * \real{0.2000}}
  >{\raggedright\arraybackslash}p{(\linewidth - 8\tabcolsep) * \real{0.2000}}@{}}
\toprule\noalign{}
\begin{minipage}[b]{\linewidth}\raggedright
split
\end{minipage} & \begin{minipage}[b]{\linewidth}\raggedright
collective entropy
\end{minipage} & \begin{minipage}[b]{\linewidth}\raggedright
information gain
\end{minipage} & \begin{minipage}[b]{\linewidth}\raggedright
weighted Gini
\end{minipage} & \begin{minipage}[b]{\linewidth}\raggedright
Gini gain
\end{minipage} \\
\midrule\noalign{}
\endhead
\bottomrule\noalign{}
\endlastfoot
Home owner & \blank[0.9in] & \blank[0.9in] & \blank[0.9in] &
\blank[0.9in] \\
Married & \blank[0.9in] & \blank[0.9in] & \blank[0.9in] &
\blank[0.9in] \\
\end{longtable}
}

Greedy choice: \blank[1.2in]

\answerspace{3.5in}

\begin{enumerate}
\def\labelenumi{(\alph{enumi})}
\setcounter{enumi}{2}
\item
  Grow the tree one more level: on the impure child from (b), split on
  the remaining attribute, and compute its information gain. Write out
  the resulting two-level tree (the leaf labels are the majority class
  in each leaf) and its training accuracy on these ten rows.\nopagebreak

  Gain of the second split: \blank[0.8in]\quad Training accuracy:
  \blank[0.8in]
\end{enumerate}

\answerspace{3in}

\begin{enumerate}
\def\labelenumi{(\alph{enumi})}
\setcounter{enumi}{3}
\tightlist
\item
  One leaf of your tree is still impure. Explain in two or three
  sentences why a tree grown to purity on a real dataset (with more
  attributes to split on) tends to overfit, and say what each of the
  following would do to the tree in (c): \texttt{max\_depth\ =\ 1};
  \texttt{min\_samples\_leaf\ =\ 3} (a leaf must contain at least 3
  training rows).
\end{enumerate}

\answerspace{2.2in}

\answerpage

\subsection{Problem 3: What averaging can and cannot
buy}\label{problem-3-what-averaging-can-and-cannot-buy}

\emph{Practices: reasoning about ensembles with one identity, and
separating bagging, random forests, and boosting by which error term
they attack.}

Suppose \(B\) predictors each have variance \(\sigma^2\) at some input,
and every pair has correlation \(\rho\). The variance of their average
is
\[\operatorname{Var}\!\Big(\frac{1}{B}\sum_{b=1}^{B}\hat y_b\Big) = \rho\,\sigma^2 + \frac{(1-\rho)\,\sigma^2}{B}.\]

\begin{enumerate}
\def\labelenumi{(\alph{enumi})}
\tightlist
\item
  Take \(\sigma^2 = 1\). Compute the variance of the average for
  \((\rho, B) = (0.8, 10)\), \((0.8, 100)\), and \((0.2, 100)\). What
  does the variance approach as \(B \to \infty\) at each \(\rho\)? In
  one sentence: what is the more effective lever, more trees or less
  correlated trees?
\end{enumerate}

{\def\LTcaptype{none} % do not increment counter
\begin{longtable}[]{@{}
  >{\raggedright\arraybackslash}p{(\linewidth - 10\tabcolsep) * \real{0.1667}}
  >{\raggedright\arraybackslash}p{(\linewidth - 10\tabcolsep) * \real{0.1667}}
  >{\raggedright\arraybackslash}p{(\linewidth - 10\tabcolsep) * \real{0.1667}}
  >{\raggedright\arraybackslash}p{(\linewidth - 10\tabcolsep) * \real{0.1667}}
  >{\raggedright\arraybackslash}p{(\linewidth - 10\tabcolsep) * \real{0.1667}}
  >{\raggedright\arraybackslash}p{(\linewidth - 10\tabcolsep) * \real{0.1667}}@{}}
\toprule\noalign{}
\begin{minipage}[b]{\linewidth}\raggedright
\((\rho, B)\)
\end{minipage} & \begin{minipage}[b]{\linewidth}\raggedright
\((0.8, 10)\)
\end{minipage} & \begin{minipage}[b]{\linewidth}\raggedright
\((0.8, 100)\)
\end{minipage} & \begin{minipage}[b]{\linewidth}\raggedright
\((0.2, 100)\)
\end{minipage} & \begin{minipage}[b]{\linewidth}\raggedright
limit, \(\rho = 0.8\)
\end{minipage} & \begin{minipage}[b]{\linewidth}\raggedright
limit, \(\rho = 0.2\)
\end{minipage} \\
\midrule\noalign{}
\endhead
\bottomrule\noalign{}
\endlastfoot
variance of the average & \blank[0.8in] & \blank[0.8in] & \blank[0.8in]
& \blank[0.8in] & \blank[0.8in] \\
\end{longtable}
}

\answerspace{1.5in}

\begin{enumerate}
\def\labelenumi{(\alph{enumi})}
\setcounter{enumi}{1}
\tightlist
\item
  Verify the identity for \(B = 2\) using
  \(\operatorname{Var}(X+Y) = \operatorname{Var}X + \operatorname{Var}Y + 2\operatorname{Cov}(X,Y)\),
  \(\operatorname{Var}(aX) = a^2\operatorname{Var}X\), and
  \(\operatorname{Cov}(X,Y) = \rho\,\sigma^2\) when both variances equal
  \(\sigma^2\).
\end{enumerate}

\answerspace{2in}
\answerpage

\begin{enumerate}
\def\labelenumi{(\alph{enumi})}
\setcounter{enumi}{2}
\tightlist
\item
  \emph{Bagging} grows each tree on a bootstrap sample of the rows and
  averages; a \emph{random forest} does the same but, at every split,
  considers only a random subset of \texttt{max\_features} of the
  attributes. In terms of \(\rho\), \(\sigma^2\), and \(B\), say what
  each ingredient buys, and why the L08 activity found that bagging
  alone (all features at every split) barely improved on one deep tree
  when a few strong attributes dominate. What might a very small
  \texttt{max\_features} cost you?
\end{enumerate}

\answerspace{2.2in}

\begin{enumerate}
\def\labelenumi{(\alph{enumi})}
\setcounter{enumi}{3}
\tightlist
\item
  \emph{Boosting} (not covered in lecture --- here is all you need) fits
  many \emph{shallow} trees in sequence, each one trained on the errors
  the ensemble has made so far, and adds them up. Which of bias and
  variance is a shallow tree short on, so which term is boosting mainly
  reducing? Contrast with the forest in one sentence: why deep trees for
  a forest but shallow trees for boosting?
\end{enumerate}

\answerspace{1.8in}




\end{document}
